\documentclass[11pt]{article}

\usepackage[utf8]{inputenc}
\usepackage[T1]{fontenc}
\usepackage{lmodern}
\usepackage{geometry}
\usepackage{amsmath,amssymb,amsfonts,amsthm,mathtools}
\usepackage{bm}
\usepackage{enumitem}
\usepackage{microtype}
\usepackage{hyperref}
\usepackage{titlesec}
\usepackage{booktabs}
\usepackage{array}
\usepackage{setspace}
\usepackage{mathrsfs}
\usepackage{physics}

\geometry{margin=1in}
\hypersetup{colorlinks=true, linkcolor=black, urlcolor=blue, citecolor=black}
\setlist[enumerate]{topsep=0.4em,itemsep=0.35em}
\setlist[itemize]{topsep=0.3em,itemsep=0.25em}
\titleformat{\section}{\large\bfseries}{\thesection.}{0.5em}{}
\titleformat{\subsection}{\normalsize\bfseries}{\thesubsection.}{0.5em}{}
\setstretch{1.06}

\newcommand{\Aether}{\AE{}ther}
\newcommand{\Aflow}{\AE{}ther-flow}
\newcommand{\TheoryName}{\Aflow{} Interpretation of Relativity}
\newcommand{\TheoryProgram}{\Aether{} / \Aflow{} framework}
\newcommand{\ObserverMetric}{g^{\mathrm{br}}}
\newcommand{\CandidateMetric}{g^{\mathrm{cand}}}
\newcommand{\CompletedMetric}{g^{\sharp}}
\newcommand{\BaseShift}{\Sigma^{\mathrm{IER}}}
\newcommand{\ResponseShift}{\Sigma^{\mathrm{resp}}}
\newcommand{\CompShift}{\Sigma^{\mathrm{cmp}}}
\newcommand{\ResidualCore}{\mathcal V_{\mu\nu}^{\mathrm{res}}}
\newcommand{\ResidualDec}{\mathcal D_{\mu\nu}^{\mathrm{res}}}
\newcommand{\ReducedEin}{\mathcal L_{\ObserverMetric}}
\newcommand{\GaugeBook}{\mathcal G_{\ObserverMetric}}
\newcommand{\EinQuad}{\mathcal Q_{\ge 2}^{\mathrm{Ein}}}
\newcommand{\CompMix}{\mathcal M_{\mu\nu}^{\mathrm{cmp}}}
\newcommand{\CompQuad}{\mathcal Q_{\mu\nu}^{\mathrm{cmp}}}
\newcommand{\DeltaTComp}{\Delta T_{\mu\nu}^{\mathrm{cmp}}}
\newcommand{\ObsBurden}{\mathfrak B_{\mu\nu}^{\mathrm{obs}}}
\newcommand{\epsIR}{\varepsilon_{\mathrm{IR}}}

\newtheorem{definition}{Definition}
\newtheorem{proposition}{Proposition}
\newtheorem{remark}{Remark}
\newtheorem{corollary}{Corollary}
\newtheorem{theorem}{Theorem}

\title{\textbf{The \TheoryName{}}\\[0.4em]
\large Substrate-Response / Self-Response Charge-Polarization Candidate\\
\large Exact-Preservation Observer-Equation Burden Sector-Classification Note}
\author{}
\date{}

\begin{document}

\maketitle

\begin{abstract}
The observer-equation burden-collapse note reduced the exact-preservation line to one explicit post-completion tensor:
\[
\ObsBurden
=
\ResidualDec+\CompMix+\CompQuad-8\pi G_N\,\DeltaTComp.
\]
The present manuscript stays on that same one-metric line and attacks \(\ObsBurden\) directly. It does not go one layer deeper. Instead it classifies the four displayed sectors first on matter-free reduced data and then on the full completed branch.

The result is sharp at the disciplined scope now required. On matter-free reduced data, \(\DeltaTComp\) is absent, while \(\ResidualDec\), \(\CompMix\), and \(\CompQuad\) are benchmark-decoupled. Consequently \(\ObsBurden\) is benchmark-decoupled there and satisfies
\[
\ObsBurden=\mathcal O(\epsIR^6)
\]
on every controlled infrared family. On the full one-metric branch, \(\ResidualDec\), \(\CompMix\), \(\CompQuad\), and \(\DeltaTComp\) are all benchmark-decoupled, so
\[
\ObsBurden=\mathcal O(\epsIR^4)
\]
and still carries no genuine observer-side remainder at benchmark-facing scope.

This is a real observer-equation gain on the same exact-preservation branch. The live burden left by the burden-collapse note is no longer merely exposed; it is now safely classified sector by sector in the sense relevant to the benchmark ``no genuine observer-side remainder'' gate. The note does not claim the stronger exact identity \(\ObsBurden=0\), and it does not claim a completed first-principles derivation of GR.
\end{abstract}

\section{Introduction}

The previous observer-equation burden-collapse note made the current priority unambiguous. Once the completed observer equation had been written explicitly on the preserved one-metric branch, deeper exact-preservation relays ceased to be the default primary move. The live benchmark-facing task became the status of one explicit tensor,
\[
\ObsBurden
=
\ResidualDec+\CompMix+\CompQuad-8\pi G_N\,\DeltaTComp.
\]
The present note performs that next step directly.

Its scope is narrow on purpose. It does not invent a new branch, a new completion solve, or a deeper primitive origin. It uses only the already-preserved one-metric exact-preservation package and asks a single question: for each displayed term in \(\ObsBurden\), is the sector absent, gauge exact, constraint exact, or benchmark-decoupled in the sense relevant to the benchmark derivation gates?

The answer is enough to change manuscript status. The note proves that the gauge-exact and constraint-exact sectors have already been removed before \(\ObsBurden\) is written, and that every term that does remain in \(\ObsBurden\) is either absent or benchmark-decoupled. This is exactly the kind of observer-equation gain that the current research rule requires.

\paragraph{Framework and claim boundary.}
\TheoryProgram{} denotes the broader \Aether{} / \Aflow{} framework built on the claims that the \Aether{} is the underlying four-dimensional substrate of reality, the \Aflow{} is its intrinsic ordered motion, observed three-dimensional space is the local experiential slice of that deeper substrate, S-time is the experienced order of change arising from matter, light, and the \Aflow{}, and observed expansion is the three-dimensional appearance of deeper four-dimensional ordered motion. Within that framework, \TheoryName{} denotes the exact-closure theory in which the effective gravitational dynamics are adopted to be exactly Einsteinian with universal matter coupling. In the active sequence, \emph{adoption} means use of that established relativistic dynamics without claiming substrate derivation, while \emph{derivation} is reserved for a first-principles recovery from explicit substrate variables.

The present note stays strictly inside that boundary. It does not claim that GR has now been derived from substrate microphysics. It does not claim the strict tensor identity \(\ObsBurden=0\). Its exact positive claim is narrower: on the preserved one-metric exact-preservation line, the remaining observer-equation burden has now been safely classified sector by sector at benchmark-facing scope.

\section{Fixed One-Metric Inputs}

\begin{definition}[Fixed exact-preservation observer package]
\label{def:fixed}
The present note treats the following observer-level data as fixed inputs from the same exact-preservation line.
\begin{enumerate}
    \item The candidate and completed operative metrics are
    \begin{equation}
    \CandidateMetric
    =
    \ObserverMetric+\BaseShift+\ResponseShift,
    \qquad
    \CompletedMetric
    =
    \ObserverMetric+\BaseShift+\ResponseShift+\CompShift.
    \label{eq:metrics}
    \end{equation}
    \item The same completion solve is kept fixed:
    \begin{equation}
    \ReducedEin(\CompShift)=-\ResidualCore.
    \label{eq:solve}
    \end{equation}
    \item The completed observer equation is
    \begin{equation}
    G_{\mu\nu}[\CompletedMetric]
    =
    8\pi G_N\,T_{\mu\nu}^{\mathrm{matter}}[\CompletedMetric,\psi]
    +\GaugeBook(\CompShift)_{\mu\nu}
    +\ObsBurden,
    \label{eq:completed}
    \end{equation}
    with
    \begin{equation}
    \ObsBurden
    \equiv
    \ResidualDec+\CompMix+\CompQuad-8\pi G_N\,\DeltaTComp.
    \label{eq:burden}
    \end{equation}
    \item The same one operative metric and ordinary matter route are preserved. In particular, \(\DeltaTComp\) is only the matter re-expansion
    \begin{equation}
    \DeltaTComp
    \equiv
    T_{\mu\nu}^{\mathrm{matter}}[\CompletedMetric,\psi]
    -
    T_{\mu\nu}^{\mathrm{matter}}[\CandidateMetric,\psi].
    \label{eq:deltat}
    \end{equation}
    \item The previously recorded residual-sector verdict on the same charge-polarization line already classifies \(\ResidualDec\) as decoupled in the benchmark sense.
    \item On every controlled infrared family of the preserved Phase F package,
    \begin{equation}
    \BaseShift=\mathcal O(\epsIR^2),
    \qquad
    \ResponseShift=\mathcal O(\epsIR^4),
    \qquad
    \CompShift=\mathcal O(\epsIR^4),
    \label{eq:orders}
    \end{equation}
    and the admissible homogeneous completion mode is fixed to zero, so the completion adds no independent unsuppressed linear observer mode.
\end{enumerate}
\end{definition}

\begin{definition}[Benchmark-decoupled sector]
\label{def:benchdec}
An observer-side tensor sector \(X_{\mu\nu}\) on the fixed branch of Definition \ref{def:fixed} is \emph{benchmark-decoupled} if all of the following hold.
\begin{enumerate}
    \item \(X_{\mu\nu}\) is generated only by already-screened response/completion variables or by one-metric matter re-expansion under the same operative metric.
    \item \(X_{\mu\nu}\) introduces no independent unsuppressed linear observer mode on the admissible homogeneous benchmark background.
    \item \(X_{\mu\nu}\) introduces no second operative metric and no second universal matter law.
    \item On every controlled infrared family, \(X_{\mu\nu}=\mathcal O(\epsIR^p)\) for some \(p>0\).
\end{enumerate}
\end{definition}

\begin{remark}
Definition \ref{def:benchdec} is exactly the present note's disciplined replacement for informal ``harmless remainder'' language. It is the sense in which a sector can fail to vanish identically while still not counting as a genuine observer-side remainder for the benchmark gate.
\end{remark}

\begin{proposition}[Finite sums preserve benchmark decoupling]
\label{prop:sum}
If \(X_{\mu\nu}\) and \(Y_{\mu\nu}\) are benchmark-decoupled on the fixed branch, then \(X_{\mu\nu}+Y_{\mu\nu}\) is benchmark-decoupled as well.
\end{proposition}

\begin{proof}
Items (1)--(3) of Definition \ref{def:benchdec} are preserved under finite sums, and item (4) is preserved by taking the smaller infrared order of the summands.
\end{proof}

\begin{definition}[Matter-free reduced data]
\label{def:matterfree}
The \emph{matter-free reduced-data regime} on the fixed exact-preservation branch is the regime in which the reduced slice package and the completion solve \eqref{eq:solve} are imposed and the ordinary matter sector is absent, so
\begin{equation}
T_{\mu\nu}^{\mathrm{matter}}[\CandidateMetric,\psi]=0,
\qquad
T_{\mu\nu}^{\mathrm{matter}}[\CompletedMetric,\psi]=0.
\label{eq:matterfree}
\end{equation}
\end{definition}

\section{What Is Already Outside \texorpdfstring{\(\ObsBurden\)}{ObsBurden}}

\begin{proposition}[Gauge-exact and constraint-exact sectors are already extracted]
\label{prop:outside}
On the fixed one-metric exact-preservation branch, the gauge-exact bookkeeping sector appears explicitly as \(\GaugeBook(\CompShift)\) in \eqref{eq:completed}, while the constraint-exact elimination sectors have already been removed before \(\ObsBurden\) is defined. Consequently the sector-by-sector classification problem for \(\ObsBurden\) reduces exactly to the four terms displayed in \eqref{eq:burden}.
\end{proposition}

\begin{proof}
Equation \eqref{eq:completed} already separates the reduced gauge-bookkeeping tensor \(\GaugeBook(\CompShift)\) from \(\ObsBurden\). The candidate-line residual verdict that feeds into Definition \ref{def:fixed} was written only after the eliminated source, polarization, and reduced-slice equations had been imposed, so the surviving constraint-exact terms are already absent from \(\ResidualDec\) and hence from \(\ObsBurden\). Therefore no additional hidden gauge-exact or constraint-exact sector remains to be classified inside \eqref{eq:burden}; the only live terms are \(\ResidualDec\), \(\CompMix\), \(\CompQuad\), and \(\DeltaTComp\).
\end{proof}

\begin{remark}
Proposition \ref{prop:outside} explains why the present note proves only absence or benchmark decoupling for the displayed burden terms. The gauge-exact and constraint-exact work has already been done earlier in the same manuscript line.
\end{remark}

\section{Matter-Free Reduced-Data Classification}

\begin{proposition}[The inherited block \(\ResidualDec\) is benchmark-decoupled on matter-free reduced data]
\label{prop:mfres}
On the regime of Definition \ref{def:matterfree}, the inherited burden term \(\ResidualDec\) is benchmark-decoupled. On every controlled infrared family in that regime,
\begin{equation}
\ResidualDec=\mathcal O(\epsIR^6).
\label{eq:mfres}
\end{equation}
\end{proposition}

\begin{proof}
Definition \ref{def:fixed} already records that \(\ResidualDec\) is decoupled on the same charge-polarization line. In the matter-free regime \eqref{eq:matterfree}, the candidate-branch matter re-expansion piece inside that inherited block vanishes, so only the already-screened response-generated nonlinear pieces remain. Those pieces are at least bilinear in the response data and therefore begin two reduced-data orders below the leading quartic observer scale. Hence \(\ResidualDec\) satisfies Definition \ref{def:benchdec} and obeys \eqref{eq:mfres}.
\end{proof}

\begin{proposition}[The completion-generated Einstein sectors are benchmark-decoupled on matter-free reduced data]
\label{prop:mfcomp}
On the regime of Definition \ref{def:matterfree}, both \(\CompMix\) and \(\CompQuad\) are benchmark-decoupled. More precisely,
\begin{equation}
\CompMix=\mathcal O(\epsIR^6),
\qquad
\CompQuad=\mathcal O(\epsIR^8).
\label{eq:mfcomp}
\end{equation}
\end{proposition}

\begin{proof}
By definition, \(\CompMix\) contains only nonlinear Einstein terms involving at least one factor of \(\CompShift\) and at least one factor of \(\BaseShift+\ResponseShift\). Using \eqref{eq:orders},
\[
\BaseShift+\ResponseShift=\mathcal O(\epsIR^2),
\qquad
\CompShift=\mathcal O(\epsIR^4),
\]
so every contribution to \(\CompMix\) is at least \(\mathcal O(\epsIR^6)\). Likewise \(\CompQuad\) is built entirely from terms at least quadratic in \(\CompShift\), hence \(\CompQuad=\mathcal O(\epsIR^8)\). Because the admissible homogeneous completion mode is fixed to zero, neither tensor carries an independent unsuppressed linear observer imprint. Both are generated only by the already-screened completion/response variables on the same one operative metric. Therefore both sectors satisfy Definition \ref{def:benchdec}.
\end{proof}

\begin{proposition}[The matter re-expansion sector is absent on matter-free reduced data]
\label{prop:mfmatter}
On the regime of Definition \ref{def:matterfree},
\begin{equation}
\DeltaTComp=0.
\label{eq:mfmatter}
\end{equation}
Hence the displayed \(-8\pi G_N\,\DeltaTComp\) contribution to \(\ObsBurden\) is absent there.
\end{proposition}

\begin{proof}
Equation \eqref{eq:matterfree} states that the ordinary matter stress vanishes on both \(\CandidateMetric\) and \(\CompletedMetric\). Substituting into \eqref{eq:deltat} gives \eqref{eq:mfmatter}.
\end{proof}

\begin{theorem}[Matter-free sector-by-sector classification of \(\ObsBurden\)]
\label{thm:matterfree}
On matter-free reduced data, the burden sectors are classified as follows.
\begin{enumerate}
    \item \(\ResidualDec\) is benchmark-decoupled.
    \item \(\CompMix\) is benchmark-decoupled.
    \item \(\CompQuad\) is benchmark-decoupled.
    \item \(\DeltaTComp\) is absent.
\end{enumerate}
Consequently \(\ObsBurden\) is benchmark-decoupled on matter-free reduced data, and on every controlled infrared family,
\begin{equation}
\ObsBurden=\mathcal O(\epsIR^6).
\label{eq:mforder}
\end{equation}
\end{theorem}

\begin{proof}
Items (1)--(4) are Propositions \ref{prop:mfres}, \ref{prop:mfcomp}, and \ref{prop:mfmatter}. The displayed tensor \(-8\pi G_N\,\DeltaTComp\) is absent by item (4), and the remaining three terms are benchmark-decoupled; Proposition \ref{prop:sum} then shows that their sum \(\ObsBurden\) is benchmark-decoupled. The order bound \eqref{eq:mforder} follows from \eqref{eq:mfres} and \eqref{eq:mfcomp}.
\end{proof}

\section{Full One-Metric Branch Classification}

\begin{proposition}[The inherited block \(\ResidualDec\) remains benchmark-decoupled on the full branch]
\label{prop:fullres}
On the full one-metric exact-preservation branch, \(\ResidualDec\) is benchmark-decoupled.
\end{proposition}

\begin{proof}
This is one of the fixed inputs in Definition \ref{def:fixed}. The structural reason is unchanged after the completion step: \(\ResidualDec\) remains the already-recorded response-generated / one-metric-re-expansion block, not a new observer-accessible law. It therefore satisfies Definition \ref{def:benchdec}.
\end{proof}

\begin{proposition}[The completion-generated Einstein sectors remain benchmark-decoupled on the full branch]
\label{prop:fullcomp}
On the full one-metric branch, \(\CompMix\) and \(\CompQuad\) are benchmark-decoupled. On every controlled infrared family,
\begin{equation}
\CompMix=\mathcal O(\epsIR^6),
\qquad
\CompQuad=\mathcal O(\epsIR^8).
\label{eq:fullcomp}
\end{equation}
\end{proposition}

\begin{proof}
The same order calculation as in Proposition \ref{prop:mfcomp} uses only the fixed branch orders \eqref{eq:orders}, not the matter-free specialization. The no-independent-linear-mode statement is also unchanged, because the homogeneous completion mode is still fixed to zero. Since \(\CompMix\) and \(\CompQuad\) are purely completion-generated nonlinear Einstein sectors on the same operative metric, they again satisfy Definition \ref{def:benchdec}.
\end{proof}

\begin{proposition}[The matter re-expansion sector is benchmark-decoupled on the full branch]
\label{prop:fullmatter}
On the full one-metric exact-preservation branch, \(\DeltaTComp\) is benchmark-decoupled. On every controlled infrared family,
\begin{equation}
\DeltaTComp=\mathcal O(\epsIR^4).
\label{eq:fullmatter}
\end{equation}
\end{proposition}

\begin{proof}
By Definition \ref{def:fixed}, \(\DeltaTComp\) is not a second source law. It is only the re-expansion of ordinary matter stress under the same single operative metric \(\CompletedMetric\). Because the completion adds no independent linear observer mode, this re-expansion introduces no new unsuppressed linear observer sector. Smooth dependence of ordinary matter stress on the operative metric then combines with \(\CompShift=\mathcal O(\epsIR^4)\) from \eqref{eq:orders} to give \eqref{eq:fullmatter}. Therefore \(\DeltaTComp\) satisfies Definition \ref{def:benchdec}.
\end{proof}

\begin{theorem}[Full-branch sector-by-sector classification of \(\ObsBurden\)]
\label{thm:full}
On the full one-metric exact-preservation branch, the burden sectors are classified as follows.
\begin{enumerate}
    \item \(\ResidualDec\) is benchmark-decoupled.
    \item \(\CompMix\) is benchmark-decoupled.
    \item \(\CompQuad\) is benchmark-decoupled.
    \item \(\DeltaTComp\) is benchmark-decoupled.
\end{enumerate}
Consequently \(\ObsBurden\) is benchmark-decoupled on the full completed branch, and on every controlled infrared family,
\begin{equation}
\ObsBurden=\mathcal O(\epsIR^4).
\label{eq:fullorder}
\end{equation}
\end{theorem}

\begin{proof}
Items (1)--(4) are Propositions \ref{prop:fullres}, \ref{prop:fullcomp}, and \ref{prop:fullmatter}. Proposition \ref{prop:sum} therefore implies that the sum \eqref{eq:burden} is benchmark-decoupled. The order bound \eqref{eq:fullorder} follows because \(\ResidualDec=\mathcal O(\epsIR^4)\) is already part of the preserved controlled-limit package, while \eqref{eq:fullcomp} and \eqref{eq:fullmatter} give higher or equal order for the other displayed sectors.
\end{proof}

\begin{remark}
Theorem \ref{thm:full} is the key observer-equation gain of the present note. The burden left by the previous note is no longer an unclassified post-completion tensor. Every displayed sector has now been safely placed into the benchmark taxonomy required by the derivation gates.
\end{remark}

\section{Gate Consequence}

\begin{theorem}[Benchmark gate consequence on the same one-metric line]
\label{thm:gate}
At benchmark-facing scope on the preserved one-metric exact-preservation branch, no genuine observer-side remainder remains inside \(\ObsBurden\). More precisely:
\begin{enumerate}
    \item on matter-free reduced data, every displayed burden term is either absent or benchmark-decoupled by Theorem \ref{thm:matterfree};
    \item on the full one-metric branch, every displayed burden term is benchmark-decoupled by Theorem \ref{thm:full}.
\end{enumerate}
Therefore the live burden isolated by the observer-equation burden-collapse note has now been safely classified in the sense relevant to the benchmark ``no genuine observer-side remainder'' gate.
\end{theorem}

\begin{proof}
Theorems \ref{thm:matterfree} and \ref{thm:full} give the complete sector-by-sector classification of all four displayed terms in \eqref{eq:burden}. Proposition \ref{prop:outside} shows that no hidden gauge-exact or constraint-exact sector remains inside \(\ObsBurden\). Hence the benchmark classification problem left open by the burden-collapse note is resolved at the disciplined scope stated here.
\end{proof}

\begin{corollary}[What this gain does and does not claim]
\label{cor:scope}
The present note is a genuine Gate (5) gain on the same one-metric observer-equation line, but it does not yet prove the stronger exact identity
\[
\ObsBurden=0.
\]
Accordingly:
\begin{enumerate}
    \item the note safely classifies the remaining observer-side burden at benchmark-facing scope;
    \item it does not claim a completed first-principles derivation of GR;
    \item any later primary note on this line must now either sharpen this classification toward exact absence/identity where honest, or else prove a genuine necessity, uniqueness, or compression theorem for the preserved exact-preservation package.
\end{enumerate}
\end{corollary}

\begin{proof}
The Gate (5) language concerns the absence of any \emph{genuine} observer-side remainder, and Theorem \ref{thm:gate} gives exactly the required benchmark-safe classification. The stricter identity \(\ObsBurden=0\) would be an exact Einstein-closure statement beyond what has been proved here. The final item follows because the sector-classification burden exposed by the previous note has now been handled.
\end{proof}

\section{Claim-Boundary Verdict}

\begin{theorem}[Observer-burden sector-classification verdict]
\label{thm:verdict}
At the disciplined scope of the present note, the exact positive result is the following.
\begin{enumerate}
    \item The same candidate metric, the same completed metric, the same completion solve, and the same one operative metric are kept fixed.
    \item Gauge-exact bookkeeping remains isolated in \(\GaugeBook(\CompShift)\), and the constraint-exact elimination sectors are already outside \(\ObsBurden\).
    \item On matter-free reduced data, \(\ResidualDec\), \(\CompMix\), and \(\CompQuad\) are benchmark-decoupled while \(\DeltaTComp\) is absent, so \(\ObsBurden\) is benchmark-decoupled with \(\ObsBurden=\mathcal O(\epsIR^6)\).
    \item On the full one-metric branch, \(\ResidualDec\), \(\CompMix\), \(\CompQuad\), and \(\DeltaTComp\) are benchmark-decoupled, so \(\ObsBurden\) is benchmark-decoupled with \(\ObsBurden=\mathcal O(\epsIR^4)\).
    \item Therefore the burden isolated by the previous observer-equation note is no longer an unclassified genuine observer-side remainder at benchmark-facing scope.
    \item This is a real observer-equation gain on the same one-metric line and a genuine advance of the benchmark no-genuine-remainder gate.
    \item Stronger promotion language is still not justified here, because the note does not prove the strict exact identity \(\ObsBurden=0\) and does not claim a completed first-principles substrate derivation of GR.
\end{enumerate}
\end{theorem}

\begin{proof}
Item (1) is Definition \ref{def:fixed}. Item (2) is Proposition \ref{prop:outside}. Item (3) is Theorem \ref{thm:matterfree}. Item (4) is Theorem \ref{thm:full}. Item (5) is Theorem \ref{thm:gate}. Item (6) is Theorem \ref{thm:gate} together with Corollary \ref{cor:scope}. Item (7) is exactly the claim boundary stated in the introduction and Corollary \ref{cor:scope}.
\end{proof}

\section{Conclusion}

The previous observer-equation note made the next primary move explicit: stay on the same one-metric exact-preservation line and classify \(\ObsBurden\) itself rather than defaulting to another deeper-origin relay. The present manuscript does exactly that. It keeps the same completed observer equation and decomposes the remaining burden sector by sector.

The classification is now explicit. On matter-free reduced data, the one-metric matter re-expansion term \(\DeltaTComp\) is absent, while \(\ResidualDec\), \(\CompMix\), and \(\CompQuad\) are benchmark-decoupled. On the full one-metric branch, all four displayed sectors are benchmark-decoupled. Thus the tensor left open by the burden-collapse note is no longer an unclassified genuine observer-side remainder at benchmark-facing scope.

This is the positive result recorded here. It is a real observer-equation gain and a genuine advance of the benchmark no-genuine-remainder gate on the same one-metric line. But it is still narrower than a completed first-principles derivation of GR: the stronger exact identity \(\ObsBurden=0\) has not been proved. The next honest primary move is therefore either to sharpen this benchmark-decoupled classification into a stricter exact theorem where possible, or else to prove a genuine necessity, uniqueness, or compression theorem for the preserved exact-preservation package rather than returning by default to observer-equation-neutral depth alone.

\end{document}
